% !TEX root = ../algebra_02.tex
\section{Альтернативное определение действия группы, теорема Кэли}

\begin{Statement}{Эквивалентность действия и гомоморфизма}{}
	Задать действие группы $G$ на множестве $M$ — то же самое, что задать гомоморфизм из группы $G$ в симметрическую группу $S(M)$.
\end{Statement}

\begin{proof}
	Пусть $G \curvearrowright M$. Построим отображение $\alpha \colon G \to S(M)$, где $\alpha(g) = \alpha_g$, а $\alpha_g \colon M \to M$ действует как $m \mapsto gm$.
	Заметим, что $\alpha_g^{-1} = \alpha_{g^{-1}}$, поэтому $\alpha_g$ действительно является биекцией, то есть $\alpha_g \in S(M)$.
	Проверим свойство гомоморфизма:
	\[ \alpha_{g_1 g_2}(m) = (g_1 g_2)m = g_1(g_2 m) = \alpha_{g_1}(\alpha_{g_2}(m)) \] 
	Следовательно, $\alpha_{g_1 g_2} = \alpha_{g_1} \circ \alpha_{g_2}$, и $\alpha$ — гомоморфизм.
	Обратно, если дан гомоморфизм $\alpha \colon G \to S(M)$, то действие определяется как $gm := \alpha_g(m)$, при этом нейтральный элемент переходит в $\alpha_e = \text{id}_M$.
\end{proof}

\begin{Theorem}{Теорема Кэли}{}
	Пусть $G$ — конечная группа порядка $n$. Тогда существует подгруппа $G' \le S_n$ такая, что $G \cong G'$.
\end{Theorem}

\begin{proof}
	Рассмотрим действие $G$ на себе левыми сдвигами. Оно определяет гомоморфизм $\alpha \colon G \to S(G)$, заданный как $g \mapsto (h \mapsto gh)$.
	Найдём ядро гомоморфизма:
	\[ \ker \alpha = \{g \in G \mid gh = h \quad \forall h \in G\} = \{e\} \] 
	Так как ядро тривиально, $\alpha$ является инъективным гомоморфизмом.
	Зафиксировав биекцию между множеством $G$ и $\{1, 2, \dots, n\}$, мы индуцируем изоморфизм $\lambda \colon S(G) \xrightarrow{\sim} S_n$.
	Тогда композиция отображений $\lambda \circ \alpha$ задаёт изоморфизм группы $G$ на свой образ:
	\[ G \cong \operatorname{Im}(\lambda \circ \alpha) \le S_n \] 
	Теорема доказана.
\end{proof}

\begin{Example}{Применение теоремы Кэли к четверной группе Клейна}{}
	Найдём изоморфную подгруппу $G' \le S_4$ для группы $G = (\mathbb{Z}/2\mathbb{Z}) \times (\mathbb{Z}/2\mathbb{Z})$.
	Элементы группы: $e=(0,0), a=(1,0), b=(0,1), c=(1,1)$.
	Занумеруем элементы: $1 \leftrightarrow e$, $2 \leftrightarrow a$, $3 \leftrightarrow b$, $4 \leftrightarrow c$.
	Действуя левым умножением каждого элемента на всю группу, получаем перестановки:
	\begin{itemize}
		\item Умножение на $e$ оставляет все на местах: $(1)(2)(3)(4) = \text{id}$.
		\item Умножение на $a$: $ae=a \to 2$, $aa=e \to 1$, $ab=c \to 4$, $ac=b \to 3$. Перестановка $(1\;2)(3\;4)$.
		\item Умножение на $b$: $be=b \to 3$, $ba=c \to 4$, $bb=e \to 1$, $bc=a \to 2$. Перестановка $(1\;3)(2\;4)$.
		\item Умножение на $c$: $ce=c \to 4$, $ca=b \to 3$, $cb=a \to 2$, $cc=e \to 1$. Перестановка $(1\;4)(2\;3)$.
	\end{itemize}
	Таким образом, $G \cong G' = \{\text{id}, (1\;2)(3\;4), (1\;3)(2\;4), (1\;4)(2\;3)\} \le S_4$.
\end{Example}
